\documentclass{book}

\begin{document}
  \section{Lógica Proposicional}

  \subsection{Proposiciones}

  Existen 2 tipos de proposiciones, atómicas y moleculares. Estas deben ser
  fórmulas bien formadas, es decir, solo pueden tener un valor de verdad
  asociado, y su enunciado no puede contener valores de verdad.

  \subsection{Conectores}

  En lógica, se definen conectores; estos representan un par ordenado de
  valores asociados a la relación entre 2 proposiciones atómicas. Estas
  definiciones indican qué casos deben cumplirse del enunciado, y estos
  pares ordenados se pueden representar por valores en una tabla de verdad.

  Estas definiciones nos permiten identificar algunas estrategias para
  demostrar proposiciones matemáticas. Para ello, basta probar, ya fuese
  un \textit{contraejemplo} o una construcción directa.

  \begin{definition}[title={Contraejemplo}]
  Un contraejemplo es probar que alguna de las exclusiones puede ocurrir.
  \end{definition}

  \subsection{Conjunción}

  \begin{minipage}[t]{0.48\textwidth}
  \begin{definition}[title={Conjunción}]
  El conector $p \land q$ solo permite el par ordenado $(V, V)$ y excluye
  al resto.
  \end{definition}
  \end{minipage}

  \begin{minipage}[t]{0.48\textwidth}
  \begin{tabular}{c|c|c}
  $p$ & $q$ & $p \land q$ \\
  \hline
  V & V & V \\
  V & F & F \\
  F & V & F \\
  F & F & F \\
  \end{tabular}
  \end{minipage}

  De este modo, podemos probar la conjunción como verdad por una prueba
  directa. Para probar que es falsa, basta probar un contraejemplo:

  \begin{align}
  \sim p &\to \sim(p \land q) \\
  \sim q &\to \sim(p \land q)
  \end{align}

  \subsection{Disyunción}

  \begin{minipage}[t]{0.48\textwidth}
  \begin{definition}[title={Disyunción}]
  Para el conector $p \lor q$ únicamente se excluye $(F, F)$.
  \end{definition}
  \end{minipage}

  \begin{minipage}[t]{0.48\textwidth}
  \begin{tabular}{c|c|c}
  $p$ & $q$ & $p \lor q$ \\
  \hline
  V & V & V \\
  V & F & V \\
  F & V & V \\
  F & F & F \\
  \end{tabular}
  \end{minipage}

  Dado que la lógica no funciona por enumeración, la forma correcta es
  por exclusión del contraejemplo, así basta probar que ocurre:

  \begin{equation}
  \sim p \land \sim q \equiv (F, F)
  \end{equation}

  \subsection{Implicación}

  \begin{minipage}[t]{0.48\textwidth}
  \begin{definition}[title={Implicación}]
  En lógica proposicional no se busca describir el mundo físico, es decir,
  la lógica no entiende causalidad ni temporalidad, sino que son un
  conjunto de reglas. Este conector no explica que si ocurre $p$ causa $q$,
  no es que para que ocurra $q$ deba ocurrir $p$. No se evalúa de forma
  secuencial, sino que cada uno tiene su valor. La única exclusión que
  tiene es el caso $(V, F)$.
  \end{definition}
  \end{minipage}

  \begin{minipage}[t]{0.48\textwidth}
  \begin{tabular}{c|c|c}
  $p$ & $q$ & $p \to q$ \\
  \hline
  V & V & V \\
  V & F & F \\
  F & V & V \\
  F & F & V \\
  \end{tabular}
  \end{minipage}

  Dicho de otra forma, si $p$ es falso, significa que la hipótesis no se
  cumple, por tanto no hace falta evaluar el valor de $q$. La regla de
  exclusión dice:

  \alert{No existe un caso en que $p$ sea $V$ y $q$ sea $F$}

  \alert{Si no hay promesa, no hay incumplimiento}

  El caso $(F, F)$ puede parecer contraintuitivo; ambas falsas debería ser
  falsa la proposición en general, pero esto no es así porque no se busca
  reflejar o explicar una situación del mundo físico. Esta proposición
  también se puede abordar como:

  \begin{equation}
  \sim(p \land \sim q) \equiv p \to q
  \end{equation}

  Para demostrar la identidad de dos proposiciones basta con ver sus
  valores de verdad.

  \begin{tabular}{c|c|c|c|c|c}
  $p$ & $q$ & $p \to q$ & $\sim q$ & $p \land \sim q$ & $\sim(p \land \sim q)$ \\
  \hline
  V & V & V & F & F & V \\
  V & F & F & V & V & F \\
  F & V & V & F & F & V \\
  F & F & V & V & F & V \\
  \end{tabular}

  \subsection{Bicondicional}

  \begin{minipage}[t]{0.48\textwidth}
  \begin{definition}[title={Bicondicional}]
  $p \leftrightarrow q$ es verdad cuando $p$ y $q$ tienen el mismo valor
  de verdad en todas las interpretaciones, es decir:
  \begin{equation}
  p \leftrightarrow q \equiv (p \to q) \land (q \to p)
  \end{equation}
  \end{definition}
  \end{minipage}

  \begin{minipage}[t]{0.48\textwidth}
  \begin{tabular}{c|c|c|c|c}
  $p$ & $q$ & $p \to q$ & $q \to p$ & $(p \to q) \land (q \to p)$ \\
  \hline
  V & V & V & V & V \\
  F & V & V & F & F \\
  V & F & F & V & F \\
  F & F & V & V & V \\
  \end{tabular}
  \end{minipage}

  De forma que es verdadera cuando coinciden $p$ y $q$, así:

  \begin{equation}
  p \leftrightarrow q \equiv (p \land q) \lor (\sim p \land \sim q)
  \end{equation}

  Por ello, para probar su validez se suele probar $p \to q$ y posterior
  $q \to p$, pues los valores excluidos son $(V,F)$ y $(F,V)$ para cada
  implicación, de modo que $p \leftrightarrow q$ excluye ambos.
  Formalmente, probar que su caso de exclusión no puede ocurrir se deben
  mostrar ambos.

  De modo que los valores de exclusión son aquellos en los que el conector
  es falso por definición semántica.

  \subsection{Ejemplo: Triángulo Equilátero}

  \begin{proposition}[title={Ángulos de 60° implica equilátero}]
  Si el triángulo $ABC$ tiene 2 ángulos de $60°$ entonces es equilátero.
  \end{proposition}

  \begin{proof}[title={Construcción directa}]
  En un espacio euclídeo, todo triángulo tiene $180°$, de modo que el
  tercer ángulo es de $60°$ pues:
  \begin{align}
  x + 2(60) &= 180 \\
  x &= 60°
  \end{align}
  Luego, $ABC$ es equilátero. \qed
  \end{proof}

  \obs La demostración anterior muestra que tomando a $p$ como hipótesis,
  $q$ es el resultado por la restricción del modelo (el espacio euclídeo
  en este caso). De una forma indirecta se probó que $\sim(p \land \sim q)$
  es verdad, es decir, $p \to q$.

  Al probar que no puede pasar $p$ y $\sim q$ a la vez, se demuestra que
  se excluye el único caso de exclusión semántico. Por esto en una
  demostración, asumir $p$ y deducir $q$ es equivalente a probar que no
  puede ocurrir el caso $p \land \sim q$.

  \obs Si bien, en matemáticas no siempre las demostraciones usan estas
  reglas lógicas, hay equivalencias muy útiles.

  \subsection{Eliminación de la Implicación}

  \begin{definition}[title={Eliminación de la implicación}]
  \begin{equation}
  p \to q \equiv \sim p \lor q
  \end{equation}
  ya que su valor de exclusión es justo el mismo $(V, F)$, pues al
  evaluar, la disyunción ambos serían falsos.
  \end{definition}

  \subsection{Contraposición}

  \begin{minipage}[t]{0.48\textwidth}
  \begin{definition}[title={Contraposición}]
  Se parece mucho a equivalencia: si se asume que $q$ es $F$, entonces
  $p$ también debe serlo, pues la implicación excluye el caso en que $p$
  es $V$ y $q$ es $F$.
  \begin{equation}
  p \to q \equiv \sim q \to \sim p
  \end{equation}
  \end{definition}
  \end{minipage}

  \begin{minipage}[t]{0.48\textwidth}
  \begin{tabular}{c|c|c|c|c|c}
  $p$ & $q$ & $p \to q$ & $\sim p$ & $\sim q$ & $\sim q \to \sim p$ \\
  \hline
  V & V & V & F & F & V \\
  V & F & F & F & V & F \\
  F & V & V & V & F & V \\
  F & F & V & V & V & V \\
  \end{tabular}
  \end{minipage}

  Para este punto, el cómo se aborda una proposición de equivalencia es
  que en una demostración, se trata de probar que el caso de exclusión
  efectivamente no puede ocurrir, es decir, la proposición deberá ser una
  tautología para ser verdadera.

  \subsection{Leyes de Morgan}

  \begin{theorem}[title={Leyes de De Morgan}]
  \begin{align}
  \sim(p \land q) &\equiv \sim p \lor \sim q \\
  \sim(p \lor q)  &\equiv \sim p \land \sim q
  \end{align}
  \end{theorem}

  Para que no ocurra $p \land q$, al menos una debe ser $F$. Ambas deben
  ser $F$ para que $p \lor q$ no se cumpla.

  \subsection{Formas Normales}

  Algo que puede ayudarnos a verificar factibilidad es convertir una
  proposición en su forma normal; esta es una forma estandarizada de
  escribir proposiciones usando solo ciertos conectores.

  \begin{definition}[title={Forma Normal Disyuntiva (FND)}]
  Es una conjunción de disyunciones:
  \begin{equation}
  (p_1 \lor q_1) \land (p_2 \lor q_2) \land \cdots
  \end{equation}
  \end{definition}

  \begin{definition}[title={Forma Normal Conjuntiva (FNC)}]
  Es una disyunción de conjunciones:
  \begin{equation}
  (p_1 \land q_1) \lor (p_2 \land q_2) \lor \cdots
  \end{equation}
  \end{definition}

  Toda proposición puede representarse por una FN. Para convertir una
  proposición a una forma normal, se siguen estos 3 pasos:

  \begin{enumerate}
  \item Se eliminan las implicaciones.
  \item Se expanden las negaciones usando las leyes de De Morgan, con la
  finalidad de poder agrupar linealmente.
  \item Se distribuyen los conectores.
  \end{enumerate}

  Por ejemplo:
  \begin{align}
  p \to (q \land r) &\equiv \sim p \lor (q \land r) \\
  \sim p \lor (q \land r) &\equiv (\sim p \lor q) \land (\sim p \lor r) \\
  \therefore\quad p \to (q \land r) &\equiv (\sim p \lor q) \land (\sim p \lor r)
  \end{align}

  \subsection{Modos Clásicos}

  Podemos simplificar la notación asignando nombres a ciertas
  proposiciones; estos son conocidos como los 4 modos clásicos.

  \begin{definition}[title={Modus Ponens}]
  La restricción se propaga. Estructura:
  \begin{align}
  &p \to q \\
  &p \\
  \hline
  &\therefore q
  \end{align}
  \end{definition}

  \begin{definition}[title={Modus Tollens}]
  Ya que la restricción es $p \land \sim q$ y $\sim q = V$, no puede
  ocurrir que $p = V$, de modo que solo queda $\sim p$. Estructura:
  \begin{align}
  &p \to q \\
  &\sim q \\
  \hline
  &\therefore \sim p
  \end{align}
  \end{definition}

  \begin{definition}[title={Silogismo Hipotético}]
  Es básicamente una composición de restricciones. Estructura:
  \begin{align}
  &p \to q \\
  &q \to r \\
  \hline
  &\therefore p \to r
  \end{align}
  \end{definition}

  \begin{definition}[title={Silogismo Disyuntivo}]
  Un \textit{o} excluyente. Estructura:
  \begin{align}
  &p \lor q \\
  &\sim p \\
  \hline
  &\therefore q
  \end{align}
  \end{definition}

  \subsection{Argumento Válido}

  \begin{definition}[title={Argumento}]
  Un argumento tiene la forma $(P_1, P_2, \ldots) \therefore q$, o
  formalmente $(P_1 \land P_2 \land \cdots) \to q$, y diremos que es
  válido si no existe una valuación donde las premisas son verdaderas y
  la conclusión es falsa.
  \end{definition}

  Y aquí entra una aclaración muy importante: toda regla de inferencia se
  puede justificar mostrando que determinada fórmula es una tautología.
  Para ello, tome el \textit{modus ponens}:

  Para verificarlo construimos la fórmula $((p \to q) \land p) \to q$.
  Si esta expresión es una tautología, entonces la regla es válida.

  \begin{tabular}{c|c|c|c|c}
  $p$ & $q$ & $p \to q$ & $(p \to q) \land p$ & $((p \to q) \land p) \to q$ \\
  \hline
  V & V & V & V & V \\
  V & F & F & F & V \\
  F & V & V & F & V \\
  F & F & V & F & V \\
  \end{tabular}

  Esto a su vez es aplicar la tautología como patrón, de forma que su
  notación nos permite reemplazar bloques completos con la hipótesis ya
  probada. De este modo, si tenemos
  $((p \to q) \land (q \to r)) \land p) \to r$, entonces podemos
  reemplazar todo el primer bloque por $r$ sin perder verdad.

  \begin{definition}[title={Verdad bajo valuación}]
  Decimos que una fórmula es verdad bajo una valuación $v$ si al
  evaluarla con las reglas semánticas de los conectores el resultado es $V$.
  \end{definition}

  \begin{definition}[title={Argumento válido}]
  Un argumento es válido si no existe una valuación donde todas las
  premisas sean verdaderas y las conclusiones sean falsas.
  \end{definition}
\end{document}
